\documentclass[10pt,a4paper]{amsart}
\usepackage[margin=19mm]{geometry}
\usepackage{amsmath,amssymb,mathtools}
\usepackage[scr=rsfs,cal=euler]{mathalfa}
\usepackage{xfrac}
\usepackage[unicode,hidelinks,bookmarks=false]{hyperref}
\newcommand{\ie}{i.e.\ }
\newcommand{\cf}{cf.\ }
\newcommand{\resp}{respectively\ }
\newcommand{\Subsection}{\subsection}
\providecommand{\nobreakdash}{\nobreak\hbox{-}}
\setlength{\emergencystretch}{2em}
\multlinegap0pt
\usepackage{tikz}
\usepackage{amssymb}
\usepackage{enumerate}
\usepackage{xcolor}
\newtheorem{lemma}{Lemma}
\newtheorem{theorem}{Theorem}
\usepackage{cleveref}
\crefname{lemma}{Lemma}{Lemmas}
\crefname{theorem}{Theorem}{Theorems}
\newcommand{\ex}{\mathrm{ex}}
\newcommand{\hex}{\widehat{\ex}}
\title{Balanced even cycles in signed graphs:\\ Tur\'an bounds, double covers, and parity obstructions}
\author{Lujia Wang, Xiaowei Yu}
\date{Source: arXiv:2609.10975v1 (CC BY 4.0)}
\begin{document}
\maketitle

\noindent\emph{Source and licence.} Balanced even cycles in signed graphs:\\ Tur\'an bounds, double covers, and parity obstructions. Version arXiv:2609.10975v1 (CC BY 4.0), distributed by arXiv under the Creative Commons Attribution 4.0 International licence, http://creativecommons.org/licenses/by/4.0/, as recorded in the arXiv licence metadata for this exact version and shown on the abstract page https://arxiv.org/abs/2609.10975v1. Full licence notice: \texttt{licenses/arxiv-2609.10975-CC-BY-4.0.txt}.\par\smallskip

\noindent\textit{Editorial scope.} Counted unit: the article's theorem thm:signed-BFS-bound (source0.tex lines 744--750). The article's complete original proof of it is delivered with it (source0.tex lines 751--797, one \texttt{proof} environment of 47 lines). No mathematics is written, repaired or re-derived: the statement and the proof are the article's own lines, in the article's own notation, and no retained line is substituted or re-flowed.  Retained uncounted as the source-local context the proof needs: lines 694--706.  Nothing was omitted from any retained span, so the package carries no omission record.  No retained line contains a citation, so the wrapper carries no bibliography and no bibliography entry.  No retained line contains a figure environment, an image-inclusion command or drawing code, so no image is delivered and the package declares no asset dependency.  Editorial formatting only, in addition: an AMS \texttt{amsart} preamble that restates the article's own declarations and macros used by the retained lines, namely the environments \texttt{lemma}, \texttt{theorem} on separate counters, together with the standard packages those lines need.  The retained environments print in this document as Lemma 1, Theorem 1 in order of appearance, whereas the article prints the same items with the numbering of its own full document.  The complete native source of the article as distributed in the arXiv e-print for this exact version is bundled unchanged as \texttt{sources/209972/source0.tex}.
\begin{lemma}\label{lem:monochromatic-BFS-closure}
Let $T$ be a breadth-first-search tree in the underlying graph of
$\widehat G$, with levels $L_i$, and switch $\widehat G$ so that every
edge of $T$ is positive.  Suppose that one sign class contains a chorded
cycle $H$ of length $h$, either in $G[L_i]$ or in $G(L_i,L_{i+1})$,
where $i\geq1$.  Then, for some integer $1\leq m\leq i$, the signed
graph contains balanced cycles of lengths
\[
2m+2,\ 2m+4,\ \ldots,\ 2m+2\left\lfloor\frac{h-1}{2}\right\rfloor.
\]
In particular, if $k>i$ is an integer and $h\geq2k-1$, then $\widehat G$ contains
$C_{+2k}$.
\end{lemma}

\begin{theorem}\label{thm:signed-BFS-bound}
For every $k\geq3$ and $n\geq2$,
\[
\hex(n,C_{+2k})\leq
2(k-1)n^{1+1/k}+8(k-1)n+2(k-1)n^{1-1/k}.
\]
\end{theorem}

\begin{proof}
Put $c=2(k-1)$.  We first establish the layer estimates and then use them
to prove expansion.

\emph{Layer bounds.}  In any breadth-first-search tree of a
$C_{+2k}$-free signed graph, the levels satisfy
\begin{equation}\label{eq:signed-layer-bounds}
e(G[L_i])\leq c|L_i|,\qquad
e(G(L_i,L_{i+1}))\leq c(|L_i|+|L_{i+1}|)
\quad(0\leq i<k).
\end{equation}
For $i=0$, there are no edges within $L_0$, and the number between $L_0$
and $L_1$ is $|L_1|$.  Let $i\geq1$ and switch so that the tree is
positive.  Apply the same argument to either of the graphs
$F=G[L_i]$ and $F=G(L_i,L_{i+1})$.  If $e(F)>c\,v(F)$, one sign class
has more than $(k-1)v(F)$ edges and hence average degree greater than
$2k-2$.  The first ordinary graph fact, with the integer $s=2k-2>3$,
gives a chorded cycle of length at least $2k-1$ in that sign class.
\Cref{lem:monochromatic-BFS-closure} then gives a balanced $2k$-cycle,
a contradiction.  This proves both estimates.

\emph{Expansion of the levels.}  Let $x=n^{1/k}$ and suppose that a
$C_{+2k}$-free signed graph on $n$
vertices has more than $(cx+4c+c/x)n$ edges.  Repeatedly delete vertices
of current degree at most $cx+4c+c/x$.  The process cannot delete the
whole graph, since then it would remove at most $(cx+4c+c/x)n$ edges.
Take a connected component of the remaining graph, whose minimum degree
$\delta$ satisfies $\delta>cx+4c+c/x$, and choose a breadth-first-search
tree in it.  Put $a_i=|L_i|$, so $a_0=1$.

We prove that $a_i>xa_{i-1}$ for $1\leq i\leq k$.  For $i=1$, this
follows from $a_1\geq\delta>x$.  Suppose that $2\leq i\leq k$ and
$a_{i-2}<a_{i-1}/x$.  Edges join only equal or consecutive levels, so
summing degrees in $L_{i-1}$ counts each edge within that level twice and
each edge to an adjacent level once.  Using
\eqref{eq:signed-layer-bounds}, we obtain
\begin{align*}
e(G(L_{i-1},L_i))
&\geq\delta a_{i-1}-2e(G[L_{i-1}])-e(G(L_{i-2},L_{i-1}))\\
&\geq(\delta-3c)a_{i-1}-ca_{i-2}\\
&>\left(\delta-3c-\frac cx\right)a_{i-1}
>c(x+1)a_{i-1}.
\end{align*}
On the other hand, \eqref{eq:signed-layer-bounds} bounds this quantity by
$c(a_{i-1}+a_i)$.  Hence $a_i>xa_{i-1}$, completing the induction.  It
follows that $a_k>xa_{k-1}>\cdots>x^ka_0=n$, a contradiction.
\end{proof}

\end{document}
